\chapter{Difference Becomes Countable}

\section{Barely more than telling apart}

The first chapter ended with a difference and a promise: not to over-read it, to call it a position and not a number.
This chapter makes good on the number. And the thing the whole construction keeps returning to, its small standing
surprise, shows itself at once: the number costs almost nothing. Counting, it turns out, is \emph{barely} more than
telling apart --- so barely that when the machine first asks the compiler for it, the compiler's own reply is exactly
that. \emph{``That is barely counting,''} it says; and it is right, and the ``barely'' is the whole of what this
section has to explain.

In this gauge the step from the distinguishable to the countable is not a leap onto new machinery. It is a pair of
\emph{type refinements}, each laid on the one type former the first chapter built, and each nothing more than the
previous demand with a single clause added. This is where the book's second verb finally enters and does real work.
Take the refinements one at a time.

The first refinement adds a \emph{comparison}, and with it, order. The machine asks which differences it is willing to
admit into a tally, and it answers with a threshold: a difference is admissible when it stands in the right
\(\le\)-relation to a bound --- a comparison of a stimulus against a threshold, settled the way every judgment here is
settled. This is the exact tier at which order enters the construction. The ground floor had none: it could say two
things were not the same, never which was the larger. One tier up, comparison appears --- and it appears as the
\emph{funge's like-characteristic}. To funge, recall, is to place a thing in a bag of its likes; and the admissibility
threshold is precisely the likeness the bag is gathered by --- \emph{of the kind we are willing to count}: those, and
not the others, that clear the \(\le\). The tange of the first chapter selected each thing apart; this refinement
takes the ones that pass a shared comparison and bags them together. It is the distinguishable, plus one clause.
(Notice that this sharpens the first chapter's refusal rather than softening it: ``no order at the distinguishable
tier'' still stands, and is now exact --- order enters here, one tier up, and not before.)

The second refinement lays the bag out in a row. It carries both prior demands --- distinguishable, and admissible ---
and adds one genuinely new thing: an \emph{enumeration}. There now exists a way to lay the admissible things out one
after another, each with its index, so that the bag is not merely a heap but a sequence you can walk. This is what
``countable'' actually names in the construction: not the arrival of order (order came a tier below, with
admissibility), but the arrival of \emph{enumerability} --- the admissible distinctions, set in succession. And once
they are laid out, the count is simply the \emph{height} of that enumeration --- how far up the sequence runs. A stack
of admissible distinctions, indexed, and its height is the number.

The third refinement says what a natural number then \emph{is}, and the answer is the one the machine has been
building toward without adding anything. A natural number is a distinguishable thing, and then a next one distinct
from all so far. That is Peano's~\cite{peano1889} zero and successor --- but reached here by \emph{refining the single
act the machine already had}, not posited as an axiom handed down from outside. Zero is the first thing told apart
from nothing at all; the successor is nothing but ``and another, distinct from every one so far''; and a number is the
height of the stack those successors build. Read the register point plainly: nothing has been added to the world
except the \emph{willingness to keep going}. The count does not arrive as new content; it falls out of the distinction
--- a funge of tanges, the bagging of things already told apart --- and it costs only the going-on.

Here the book owes the reader one careful line, because it is easy to claim more than the code will support. It is
natural to say --- and this book will keep saying it --- that the count is the \emph{cardinality of the funge}, the
size of the bag of like things. But that phrase is the gauge's reading, not a fact carried in the construction: the
machine builds no explicit ``size of the bag'' object. What it builds is the enumeration and its height, and the
height is what is genuinely derived; ``the count is the cardinality of the funge'' is the illuminating overlay laid on
top of it --- true as a reading, and marked as one. The distinction matters because the whole book's honesty is
exactly this care: to say what the machine derives, and to say separately what the gauge reads into it.

Two last things to fix in the reader's mind before the next section. First, on numerals: a numeral has appeared in
this section for the first time in the book --- zero, a successor, an order --- and it is worth stating the rule that
will hold for the length of it. A \emph{structural} numeral is not a reading. Zero-and-successor, order, the height of
an enumeration: these are the machinery of counting, and not one of them is a magnitude the machine has measured. When
a numeral finally does carry a measurement, it will be unmistakable, and it will not arrive until the very end.
Second, on what the machine now holds: it is still \emph{not a measurement}. It is a tally --- a count of admissible
distinctions --- and a tally is no more a reading than a difference was a number. The refusal of the first chapter
carries straight through: the machine has one more thing than it did, and it declines, again, to call that thing more
than it is. What it will take to turn a tally into a reading --- and why even the reading, at the very end, comes back
as a bracket and not a bare number --- is the work of the chapters ahead.

\section{A number is a finite condition}

A tally has a comfortable habit: it can count up forever and never has to stop. But a genuine number is not a tally
that simply kept going; it has to \emph{settle} somewhere --- to be some particular number, and not merely
more-than-the-last. And here the machine meets the first limit it flatly refuses to fake. It will not hold a completed
infinity. What it holds instead is something humbler and, this book will argue, more honest: a \emph{finite condition}
--- a value it can state right now, together with a rule for getting closer. A number, in this construction, is not a
thing that has arrived. It is a procedure that has agreed where it is going.

The type-theoretic content of that sentence is exact, and it dissolves a confusion the reader may be carrying. To
close on a real number does \emph{not} mean to hold the real number. It means to inhabit an \emph{approximation type}:
a term that packages a rational you can display this instant together with a \emph{refinement rule} --- a next-step
that improves the rational as far as anyone demands. The refinement above the countable is exactly this settling tier:
a distinguishable, counted thing acquires, here, a process that closes it. And the crucial economy is that nothing
infinite is ever stored. Even the limit the process is named for is itself held as a \emph{rational} --- a finite
thing --- so the endless destination need never be written down, and in truth could not be, not as a completed object.

Sit with what the machine actually carries at any moment, because it is less than the reader expects and it is
enough. There is a rational value, which it can show you. There is a rule that will refine that value, which it can
run. A finite datum and a finite rule --- and that is the whole of the number. The machine does not store the endless
tail of digits; it stores the \emph{ability to extend}. This is the constructive stance in its plainest form: to
possess a real is to hold the recipe, not the infinity the recipe approaches --- to have the number is to have the
refinement rule. (The analyst's word for the rate at which such a rule closes is a \emph{modulus of convergence},
after Bishop~\cite{bishop2006} --- but that is a gloss laid on from outside; what the machine carries is a next-step, a
rule that yields the following approximation, not a rate. Read ``modulus'' as the illumination and ``refinement rule''
as the thing itself.)

And notice that the ledger has not changed character at all. Each refinement is still the one act the book began with:
take one more step, tell one more thing apart. The whole tower of approximations is built out of the counting the
machine already had --- nothing new is spent to climb it but more turns. The only genuinely new clause is the closing
one, and it is a modest, decisive thing: a number is the place where the crowd of approximations agrees to stop
wandering apart. Not the place they arrive --- they never arrive --- but the place past which they no longer disagree
by anything that matters. Convergence, here, is an agreement, not an arrival.

This is the first place in the book where the machine \emph{declines} a funge, and the refusal is worth naming for
what it is. A completed real would be an actual infinite bag --- the whole endless tail of the sequence gathered up
and funged into one object, all at once. The machine will not build that bag. It keeps instead a finite tange --- the
single rational it holds now --- and a rule that will hand you the next member whenever you ask for one. The infinite
bag is refused; the finite selection and its generator are kept. Where the last section performed a funge, bagging
distinguishables into a count, this one turns a funge down --- and the turning-down is not a failure of nerve. It is
the whole discipline, arriving early: the machine keeps only what it can hold and rerun, never a totality it would
have to take on faith.

For that refusal is the seed of the shape the whole book ends on. At the very last, the machine will hold the one
coupling it reaches for in exactly this way --- not as a completed magnitude it has summed, but as a finite condition
it can state and repeat: a value in hand, a rule to sharpen it, and an honest interval where the crowd has agreed to
stop. The number that closes the book is a finite condition of the same kind as the number that opens this chapter;
what stands between them is only the long refinement. A real approached and never completed, here in the second
chapter, is the same honesty as a coupling bracketed and never claimed, there at the end. The machine learned, at the
count, to keep the recipe instead of the infinity --- and it never once goes back on it.

\section{The first loop}

The first chapter built a difference; the count stacked differences into a tally; the finite condition let the tally
settle. This section does something the construction has not yet done: it closes on itself. A thing encoded and
refined can be run back down --- decoded --- and when it comes home it carries \emph{the same fact it started from}.
This is the book's first loop, its first round-trip, and it is faithful --- \emph{lossless} --- not because
faithfulness was bolted on to protect the channel, but because at the very bottom of the trip, encoding simply
\emph{is} identity. The channel did not have to be made honest. Its honesty was already the base case.

The way to see this is to look at the relation the encoding tier carries. Alongside its settling process, that tier
holds a predicate on sequences --- a relation saying when one sequence stands to another as encoding to encoded ---
and the shape of that predicate is the whole of the section. It is defined by cases, by asking what the two sequences
look like, and there are three cases worth reading.

The first case is the floor of the loop, and it is pure identity. When both sequences have run all the way down ---
exhausted, nothing left to refine --- the relation between them collapses to a single demand: their facts must be
equal. Two spent sequences stand in the encoding relation exactly when the fact one bottoms out on is the fact the
other bottoms out on. That is the whole of ``decodes to itself'': at the very bottom of every round-trip there is
nothing but a symbol equal to itself --- the difference of the first chapter handed back unchanged after its long
detour through the count and the limit. Losslessness is not a property added to the channel to make it safe; it is
what the channel turns out to \emph{be} once you run it all the way to the ground.

The second case shows how the fact survives the descent. When both sequences are still refining --- each at some rung,
with a rank and an index --- the relation asks three things at once: that the two \emph{facts} are equal, that the
rank of the one is at most the rank of the other, and that its index is strictly the deeper. Read what that
conjunction says. The fact is held \emph{equal at every rung}; only the rank and the index are permitted to advance.
So the round-trip is not a static mirror that happens to reflect faithfully; it is a refinement that \emph{preserves
the fact as it descends}. And the channel is lossless for a precise reason: the one thing it never lets vary is the
fact. Everything that advances --- the rank climbing, the index deepening --- is the machinery of getting closer, not
the content being carried. This is the very first act of the book, telling one more thing apart, now shown to
\emph{conserve} the original distinction across the whole tower of refinements. The difference is not merely
re-derivable; it is invariant.

The third case gives the loop a direction, and the direction is where the next chapter begins. The off-diagonal cases
are not symmetric. An exhausted sequence encodes into a still-refining one; a still-refining one does not encode back
into an exhausted one. The round-trip runs one way: the spent thing encodes forward into the deeper, and not the
reverse. This one-wayness is not a technicality to be tidied away --- it is the seam the third chapter opens on. For an
encoding, as the tier's own commentary observes, hands you two things at once: an origin, and a limit away from that
origin --- which is to say a \emph{direction} and a \emph{magnitude}, the very pieces a residue is made of, one of
them covariant and one contravariant. The loop that keeps the fact also, in the act of keeping it, records how far the
fact had to travel to be kept. And that distance --- the gap the round-trip could not close to zero --- is the
antagonist the whole next chapter is about.

It is worth reading all of this once more in the two verbs. The count of the first section \emph{funged} the
differences, bagging them by their shared characteristic into an enumeration. This section runs the bag back through
the encoding relation and finds that characteristic returned intact: the base case, \emph{the facts are equal}, is
exactly the statement that the like-characteristic which formed the bag is still the one you draw back out of it. To
encode-and-refine is to bag; to decode is to draw a member back; and the floor of the relation proves that the member
you draw is the member you put in. Losslessness, said in these terms, is simply this: the funge never forgets the
tange that defined it. (That the relation \emph{is} a loop, a round-trip, a lossless channel, is a reading laid over
it --- an illuminating one, and the source flags the reach of its own framing here plainly --- while the thing
actually built is the relation and its cases. The book keeps the two apart, as always.)

One last thing, planted and not yet cashed. This first loop is the first small rehearsal of the act the whole
instrument ends on. Much later the machine will turn its construction back on itself in earnest and check that two
descriptions of one thing are the same; and it will hand back its final reading not as a number but as something you
can \emph{rerun} and get the same fact from. Both of those are this section, grown up. The round-trip that keeps the
fact is why repeatability, when it comes, will be the whole of the honesty --- and why the reading at the end can be
trusted precisely because you can send it around the loop again and it comes home unchanged.
